% Motor: XeLaTeX. Sesion 15 - material docente UTP.
% Misma estructura docente que la Sesion 14.
\documentclass[aspectratio=169,10pt,c]{beamer}
% Reutiliza sin cambios la identidad visual UTP validada en la Sesion 09.
% Reutiliza sin cambios la identidad visual UTP validada en la Sesion 08.
% Reutiliza sin cambios la identidad visual UTP validada en la Sesion 05.
% La segunda ruta permite compartir los assets comunes sin duplicar binarios.
\input{../../Sesion05/slides/header.tex}
\graphicspath{{../../Sesion05/slides/assets/}}



\graphicspath{{../../Sesion05/slides/assets/}}
% Bloques de código con el mismo estilo que la Sesion 11.
\usepackage{listings}
\lstset{basicstyle=\ttfamily\small,backgroundcolor=\color{utpGrisClaro},
  frame=none,aboveskip=4pt,belowskip=4pt,showstringspaces=false,
  columns=fullflexible,keepspaces=true,breaklines=true,
  xleftmargin=5pt,xrightmargin=5pt,framesep=5pt}


\title{Integración de Servicios Cloud GCP}
\subtitle{Repaso y ruta de certificación Cloud Engineer de Google Cloud}
\date{Semana 08 - Sesión 15}
\author{Rel Guzman}

\begin{document}

{\setbeamertemplate{footline}{}\maketitle}

\begin{frame}[c]{Revisión de la sesión anterior}
\begin{block}{La sesión anterior}
\small Vimos \textbf{WSL} y la \textbf{virtualización anidada}, publicamos la landing de Estampa en \textbf{GitHub Pages} y \textbf{Vercel}, dimos permisos con \textbf{IAM}, \textbf{habilitamos} la API de Compute Engine y creamos una \textbf{VM con NGINX} en Google Skills.
\end{block}
\begin{alertblock}{Pregunta rápida}
\small Un compañero debe \textbf{crear VMs} en tu proyecto, pero no debe poder borrar el proyecto ni dar permisos a otros. ¿Le das Visualizador, Editor o Propietario?
\end{alertblock}
\end{frame}

\begin{frame}[c]
\begin{block}{Logro de la sesión}
Al finalizar, podrás repasar con seguridad lo visto en la Sesión 14 y explicar \textbf{qué es la certificación Associate Cloud Engineer}, qué temas cubre, cómo se organiza Google Cloud (\textbf{proyectos, regiones y zonas}) y cómo empezar la \textbf{ruta Cloud Engineer} en Google Skills.
\end{block}
\end{frame}

\begin{frame}{Saberes previos}
\begin{itemize}\setlength\itemsep{10pt}
  \item ¿Qué diferencia hay entre \textbf{terminar un curso} y \textbf{aprobar una certificación}?
  \item Si creas una máquina virtual para usuarios de Perú, ¿en qué país crees que estaría el centro de datos de Google más cercano?
  \item De todo lo que hicimos en el curso, ¿qué crees que haría un \textbf{ingeniero de nube} en su trabajo diario?
\end{itemize}
\end{frame}

\begin{frame}[c]{Utilidad de esta sesión}
\begin{center}
\begin{minipage}{.86\textwidth}
{\large Hoy ordenas lo aprendido y le pones un nombre: casi todo lo que practicamos (proyectos, IAM, APIs, VMs, Cloud Shell) es parte del temario de una certificación que piden las empresas. Sales con un mapa claro para seguir aprendiendo solo.}
\end{minipage}
\end{center}
\end{frame}

% ------------------------------------------------------------ Repaso
\begin{frame}[plain]
\vfill
\begin{center}
{\fontsize{29}{34}\selectfont\bfseries Repaso de la Sesión 14}\par
\vspace{2mm}
{\small WSL, sitios estáticos, IAM y servicios}\par
\vspace{4mm}
{\color{utpRojo}\rule{8.5cm}{1.2pt}}
\end{center}
\vfill
\end{frame}

\begin{frame}{Repaso: WSL y virtualización anidada}
\begin{columns}[T,onlytextwidth]
\column{.48\textwidth}
\begin{block}{WSL}
\small \textbf{Linux real dentro de Windows}, sin reiniciar ni instalar otro sistema. Se instala con \texttt{wsl -{}-install}.\par\smallskip
\textbf{WSL 2} usa un kernel de Linux verdadero dentro de una VM ligera.
\end{block}
\begin{exampleblock}{Para qué sirve de verdad}
\small Trabajar con las mismas herramientas que el servidor: Python, Node.js, Docker Desktop, \texttt{gcloud}, \texttt{ssh}, \texttt{git}.
\end{exampleblock}
\column{.48\textwidth}
\begin{block}{Virtualización anidada}
\small Una \textbf{VM dentro de otra VM}, como las matrioskas. La VM de afuera debe recibir las funciones del procesador (VT-x o AMD-V).
\end{block}
\begin{alertblock}{La conexión entre las dos}
\small WSL 2 \textbf{ya es una VM}. Si tu Windows es una VM, usar WSL 2 ahí es virtualización anidada.
\end{alertblock}
\end{columns}
\end{frame}

\begin{frame}{Repaso: sitios estáticos}
\begin{block}{Idea clave}
\small Un sitio estático es solo \textbf{archivos} (HTML, CSS, JS). El servidor no calcula nada: los entrega tal como están. Por eso se publica gratis.
\end{block}
\renewcommand{\arraystretch}{1.3}
\centering\small
\begin{tabular}{@{}p{3.4cm}p{4.6cm}p{4.6cm}@{}}
\toprule
& \textbf{GitHub Pages} & \textbf{Vercel}\\
\midrule
Cómo se publica & Settings $\rightarrow$ Pages $\rightarrow$ rama \texttt{main} & Importar repo, preset \textbf{Other}, Deploy\\
Dirección & \texttt{usuario.github.io/repo} & \texttt{proyecto.vercel.app}\\
Vista previa por rama & No & Sí\\
\bottomrule
\end{tabular}
\par\vspace{3mm}
\begin{alertblock}{Los dos errores más comunes}
\small Rutas con \texttt{/} al inicio (\texttt{/styles.css}) y nombres con otra mayúscula (\texttt{Styles.css}): la página sale sin estilos.
\end{alertblock}
\end{frame}

\begin{frame}{Repaso: IAM}
\begin{center}
\begin{tikzpicture}[font=\small,
  caja/.style={draw=utpRojo,fill=utpRojo!10,rounded corners=4pt,minimum width=3.5cm,minimum height=1.2cm,align=center}]
  \node[caja] (q) at (0,0) {\textbf{Quién}\\\footnotesize principal};
  \node[caja] (r) at (4.6,0) {\textbf{Qué puede hacer}\\\footnotesize rol};
  \node[caja] (d) at (9.2,0) {\textbf{Dónde}\\\footnotesize recurso};
  \draw[-latex,thick] (q) -- (r);
  \draw[-latex,thick] (r) -- (d);
\end{tikzpicture}
\end{center}
\begin{columns}[T,onlytextwidth]
\column{.48\textwidth}
\begin{exampleblock}{Tipos de roles}
\small
\textbf{Básicos:} Propietario, Editor, Visualizador.\par\smallskip
\textbf{Predefinidos:} para un servicio, por ejemplo \emph{Administrador de instancias de Compute}.\par\smallskip
\textbf{Personalizados:} los armas tú.
\end{exampleblock}
\column{.48\textwidth}
\begin{block}{Mínimo privilegio}
\small Da \textbf{solo lo necesario}. Quien solo mira, Visualizador. Quien administra VMs, un rol predefinido de Compute, no Propietario.
\end{block}
\begin{alertblock}{Dónde se hace}
\small \textbf{IAM y administración} $\rightarrow$ \textbf{IAM} $\rightarrow$ \textbf{Otorgar acceso}.
\end{alertblock}
\end{columns}
\end{frame}

\begin{frame}[fragile]{Repaso: servicios y tu primera VM}
\begin{columns}[T,onlytextwidth]
\column{.48\textwidth}
\begin{block}{Habilitar servicios}
\small En un proyecto nuevo, las APIs están \textbf{apagadas}. Se encienden en \textbf{APIs y servicios} $\rightarrow$ \textbf{Biblioteca}, o con:
\begin{lstlisting}
gcloud services enable \
  compute.googleapis.com
\end{lstlisting}
\end{block}
\begin{alertblock}{IAM + servicio}
\small IAM decide \textbf{quién} lo usa; habilitar decide \textbf{si existe} en tu proyecto. Necesitas las dos.
\end{alertblock}
\column{.48\textwidth}
\begin{exampleblock}{El laboratorio de la VM}
\small
\textbf{1.} VM \texttt{gcelab} desde la consola con \emph{Permitir tráfico HTTP}.\par\smallskip
\textbf{2.} SSH, instalar \textbf{NGINX} y abrir la IP externa con \texttt{http://}.\par\smallskip
\textbf{3.} VM \texttt{gcelab2} con \texttt{gcloud compute instances create}.
\end{exampleblock}
\begin{block}{Lo que era}
\small Compute Engine es \textbf{IaaS}: te dan la máquina y tú instalas y cuidas todo lo de adentro.
\end{block}
\end{columns}
\end{frame}

% ------------------------------------------------- Ruta Cloud Engineer
\begin{frame}[plain]
\vfill
\begin{center}
{\fontsize{29}{34}\selectfont\bfseries Ruta Cloud Engineer}\par
\vspace{2mm}
{\small La certificación Associate Cloud Engineer de Google Cloud}\par
\vspace{4mm}
{\color{utpRojo}\rule{8.5cm}{1.2pt}}
\end{center}
\vfill
\end{frame}

\begin{frame}{¿Qué es una certificación de Google Cloud?}
\begin{block}{Definición}
\small Un \textbf{examen oficial} que demuestra que sabes hacer algo en Google Cloud. Un curso te enseña; la certificación \textbf{comprueba} que lo aprendiste.
\end{block}
\begin{columns}[c,onlytextwidth]
\column{.54\textwidth}
\centering
\begin{tikzpicture}[font=\footnotesize,
  esc/.style={draw=utpRojo,rounded corners=3pt,minimum width=2.1cm,align=center,anchor=south west}]
  \node[esc,fill=utpRojo!10,minimum height=1.0cm] at (0,0) {\textbf{Fundamentos}\\Cloud Digital\\Leader};
  \node[esc,fill=utpRojo!25,minimum height=1.8cm] at (2.3,0) {\textbf{Asociado}\\Associate\\Cloud Engineer};
  \node[esc,fill=utpRojo!45,minimum height=2.6cm] at (4.6,0) {\textbf{Profesional}\\Cloud Architect,\\Cloud Developer\\y otras};
  \draw[-latex,thick,utpRojo] (0.2,3.0) -- (6.5,3.0) node[midway,above,font=\scriptsize,text=black] {más experiencia};
\end{tikzpicture}
\column{.42\textwidth}
\begin{exampleblock}{Una comparación}
\small Como el \textbf{brevete}: puedes saber manejar, pero el brevete es la prueba oficial que te pide el empleador.
\end{exampleblock}
\begin{alertblock}{Hoy nos enfocamos en}
\small \textbf{Associate Cloud Engineer (ACE)}: el primer nivel técnico, pensado para quien crea y administra cosas en Google Cloud.
\end{alertblock}
\end{columns}
\end{frame}

\begin{frame}{¿Qué hace un Cloud Engineer?}
\begin{block}{En una frase}
\small \textbf{Crea, configura y cuida} lo que una empresa tiene en la nube: máquinas, bases de datos, redes, permisos y costos.
\end{block}
\begin{columns}[T,onlytextwidth]
\column{.48\textwidth}
\begin{exampleblock}{Una comparación}
\small Es el \textbf{administrador de un edificio}: no diseñó el edificio (eso es el arquitecto), pero reparte las llaves, revisa que funcione la luz, atiende fallas y controla los gastos.
\end{exampleblock}
\column{.48\textwidth}
\begin{block}{Un día normal}
\small
Crear una VM o desplegar en Cloud Run.\par\smallskip
Dar acceso a un nuevo compañero (IAM).\par\smallskip
Revisar por qué una app se cayó (registros).\par\smallskip
Poner una alerta si el gasto sube.
\end{block}
\end{columns}
\par\vspace{1mm}
\begin{alertblock}{Ya lo empezaste}
\small Crear un proyecto, habilitar una API, dar un rol y crear una VM: eso es trabajo de Cloud Engineer.
\end{alertblock}
\end{frame}

\begin{frame}{El examen Associate Cloud Engineer}
\begin{columns}[T,onlytextwidth]
\column{.48\textwidth}
\begin{block}{Cómo es}
\small
\textbf{Duración:} 2 horas.\par\smallskip
\textbf{Preguntas:} entre 50 y 60, de opción múltiple (una o varias respuestas).\par\smallskip
\textbf{Dónde:} en línea desde tu casa (con supervisión) o en un centro de evaluación.\par\smallskip
\textbf{Costo:} alrededor de 125 dólares.\par\smallskip
\textbf{Vigencia:} 3 años.
\end{block}
\column{.48\textwidth}
\begin{exampleblock}{Cómo son las preguntas}
\small Casi siempre son \textbf{situaciones}: ``Tu equipo necesita\ldots{} ¿qué haces?''. No piden memorizar, piden \textbf{elegir la mejor opción}.
\end{exampleblock}
\begin{block}{Experiencia sugerida}
\small Google recomienda unos \textbf{6 meses} de práctica con Google Cloud. Los labs cuentan como práctica.
\end{block}
\end{columns}
\par\vspace{1mm}
\begin{alertblock}{Antes de inscribirte}
\small Los datos del examen cambian: revisa la página oficial en \texttt{cloud.google.com/learn/certification/cloud-engineer}.
\end{alertblock}
\end{frame}

\begin{frame}{Los 4 bloques del examen}
\renewcommand{\arraystretch}{1.35}
\centering\small
\begin{tabular}{@{}p{0.6cm}p{5.2cm}p{6.8cm}@{}}
\toprule
& \textbf{Bloque} & \textbf{En palabras simples}\\
\midrule
1 & Preparar el entorno & Crear proyectos, vincular facturación, habilitar APIs, instalar \texttt{gcloud}\\
2 & Planificar e implementar & Elegir y crear VMs, contenedores, almacenamiento, bases de datos y redes\\
3 & Operar la solución & Mantener todo funcionando: cambios, copias, registros y monitoreo\\
4 & Acceso y seguridad & IAM, roles, cuentas de servicio y mínimo privilegio\\
\bottomrule
\end{tabular}
\par\vspace{4mm}
\begin{block}{Idea clave}
\small El examen es \textbf{práctico}: pregunta qué harías en la consola o con \texttt{gcloud}, no definiciones de memoria.
\end{block}
\end{frame}

\begin{frame}{Lo que ya viste en el curso, en el temario}
\renewcommand{\arraystretch}{1.35}
\centering\small
\begin{tabular}{@{}p{4.4cm}p{8.2cm}@{}}
\toprule
\textbf{Bloque del examen} & \textbf{Lo que ya practicaste}\\
\midrule
1. Preparar el entorno & Proyecto, crédito y facturación; \textbf{habilitar la API} de Compute Engine; Cloud Shell\\
2. Planificar e implementar & VM \texttt{gcelab} con NGINX; despliegue en PaaS (Vercel, Cloud Run)\\
3. Operar la solución & Revisar por qué la IP no abría; costos de transferencia y cómputo\\
4. Acceso y seguridad & \textbf{IAM}: roles básicos y predefinidos, mínimo privilegio\\
\bottomrule
\end{tabular}
\par\vspace{4mm}
\begin{exampleblock}{Mensaje para ti}
\small No empiezas de cero: ya tocaste los \textbf{cuatro bloques}, al menos en lo básico.
\end{exampleblock}
\end{frame}

\begin{frame}{La ruta Cloud Engineer en Google Skills}
\begin{block}{Qué es}
\small Un \textbf{camino de aprendizaje} (\emph{learning path}) en \texttt{skills.google} que ordena cursos, laboratorios e insignias para prepararte para el examen.
\end{block}
\begin{center}
\begin{tikzpicture}[font=\small,
  paso/.style={draw=utpRojo,fill=utpRojo!10,rounded corners=4pt,minimum width=3.3cm,minimum height=1.3cm,align=center}]
  \node[paso] (a) at (0,0) {\textbf{1. Aprender}\\\footnotesize cursos con videos};
  \node[paso] (b) at (4.3,0) {\textbf{2. Practicar}\\\footnotesize laboratorios reales};
  \node[paso] (c) at (8.6,0) {\textbf{3. Demostrar}\\\footnotesize insignias y examen};
  \draw[-latex,thick] (a) -- (b);
  \draw[-latex,thick] (b) -- (c);
\end{tikzpicture}
\end{center}
\begin{columns}[T,onlytextwidth]
\column{.48\textwidth}
\begin{exampleblock}{Cursos de la ruta (ejemplos)}
\small \emph{Google Cloud Fundamentals: Core Infrastructure} y la serie \emph{Essential Google Cloud Infrastructure}.
\end{exampleblock}
\column{.48\textwidth}
\begin{block}{Insignia de habilidad (\emph{skill badge})}
\small Un reto final sin instrucciones paso a paso. Si lo resuelves, ganas una insignia que puedes mostrar en LinkedIn.
\end{block}
\end{columns}
\end{frame}

% ------------------------------------------------ Conceptos basicos
\begin{frame}[plain]
\vfill
\begin{center}
{\fontsize{29}{34}\selectfont\bfseries Los primeros conceptos de la ruta}\par
\vspace{2mm}
{\small Lo básico que aparece en todos los cursos}\par
\vspace{4mm}
{\color{utpRojo}\rule{8.5cm}{1.2pt}}
\end{center}
\vfill
\end{frame}

\begin{frame}{Cómo se organiza todo: la jerarquía de recursos}
\begin{columns}[c,onlytextwidth]
\column{.46\textwidth}
\centering
\begin{tikzpicture}[font=\small,
  niv/.style={draw=utpRojo,rounded corners=4pt,minimum width=4.6cm,minimum height=0.8cm,align=center}]
  \node[niv,fill=utpRojo!45] (o) at (0,3.3) {\textbf{Organización}\\\footnotesize la empresa (\texttt{utp.edu.pe})};
  \node[niv,fill=utpRojo!30] (f) at (0,2.2) {\textbf{Carpetas}\\\footnotesize áreas: Ventas, TI};
  \node[niv,fill=utpRojo!15] (p) at (0,1.1) {\textbf{Proyectos}\\\footnotesize \texttt{mi-tienda}, \texttt{pruebas}};
  \node[niv,fill=white] (r) at (0,0) {\textbf{Recursos}\\\footnotesize VMs, buckets, bases de datos};
  \draw[-latex,thick] (o) -- (f);
  \draw[-latex,thick] (f) -- (p);
  \draw[-latex,thick] (p) -- (r);
\end{tikzpicture}
\column{.50\textwidth}
\begin{block}{El proyecto es lo más importante}
\small Todo recurso vive dentro de un \textbf{proyecto}. El proyecto tiene su propia \textbf{facturación}, sus \textbf{APIs habilitadas} y sus \textbf{permisos de IAM}.
\end{block}
\begin{exampleblock}{Una comparación}
\small Organización = el edificio; carpetas = los pisos; proyectos = los departamentos; recursos = los muebles.
\end{exampleblock}
\begin{alertblock}{Los permisos bajan}
\small Un rol dado en una carpeta \textbf{se hereda} en todos sus proyectos.
\end{alertblock}
\end{columns}
\end{frame}

\begin{frame}{Regiones y zonas: dónde vive tu recurso}
\begin{block}{Definiciones}
\small Una \textbf{región} es un lugar del mundo con centros de datos de Google (por ejemplo, \texttt{southamerica-west1} en Santiago de Chile). Cada región tiene varias \textbf{zonas}: edificios separados dentro de ella (\texttt{southamerica-west1-a}, \texttt{-b}, \texttt{-c}).
\end{block}
\begin{columns}[T,onlytextwidth]
\column{.48\textwidth}
\begin{exampleblock}{Una comparación}
\small La región es la \textbf{ciudad}; las zonas son \textbf{distritos} con su propia luz. Si se corta en uno, los otros siguen.
\end{exampleblock}
\begin{block}{Recursos zonales y regionales}
\small Una VM vive en \textbf{una zona}. Otros servicios, como Cloud Run o un bucket regional, viven en \textbf{una región}.
\end{block}
\column{.48\textwidth}
\begin{block}{Cómo elegir}
\small
\textbf{Cerca de tus usuarios:} responde más rápido.\par\smallskip
\textbf{Precio:} cambia según la región.\par\smallskip
\textbf{Plan gratuito:} algunas ofertas gratis solo aplican en ciertas regiones de EE.\,UU.
\end{block}
\begin{alertblock}{Error típico}
\small ``\texttt{gcloud} me pide zona'': falta \texttt{-{}-zone}. Recuerda el lab de la VM.
\end{alertblock}
\end{columns}
\end{frame}

\begin{frame}[fragile]{Tres formas de trabajar en Google Cloud}
\begin{columns}[T,onlytextwidth]
\column{.30\textwidth}
\begin{block}{Consola}
\small La página web \texttt{console.cloud.google.com}. Con clics. Ideal para empezar y para \textbf{ver} lo que existe.
\end{block}
\column{.30\textwidth}
\begin{block}{Cloud Shell}
\small Una terminal Linux \textbf{en el navegador}, con \texttt{gcloud} ya instalado. Gratis y sin instalar nada.
\end{block}
\column{.34\textwidth}
\begin{block}{\texttt{gcloud}}
\small La herramienta de \textbf{línea de comandos}. Se puede instalar también en tu laptop o en WSL.
\end{block}
\end{columns}
\par\vspace{1mm}
\begin{exampleblock}{Primeros comandos que verás en la ruta}
\begin{lstlisting}[basicstyle=\ttfamily\footnotesize]
gcloud auth list                                  # con que cuenta estoy
gcloud config list project                        # en que proyecto estoy
gcloud config set compute/region us-central1      # region por defecto
gcloud compute instances list                     # mis VMs
\end{lstlisting}
\end{exampleblock}
\begin{alertblock}{Consejo}
\small El examen pregunta por los dos caminos. Si lo haces con clics, busca también el comando.
\end{alertblock}
\end{frame}

\begin{frame}{Mapa de servicios para principiantes}
\renewcommand{\arraystretch}{1.2}
\centering\small
\begin{tabular}{@{}p{4.2cm}p{4.2cm}p{4.4cm}@{}}
\toprule
\textbf{Necesito\ldots} & \textbf{Servicio} & \textbf{Tipo}\\
\midrule
Una máquina completa & Compute Engine & IaaS (tú cuidas todo)\\
Correr un contenedor & Cloud Run & Sin servidor (\emph{serverless})\\
Muchos contenedores & GKE (Kubernetes) & Contenedores administrados\\
Guardar archivos e imágenes & Cloud Storage & Almacenamiento de objetos\\
Base de datos con tablas & Cloud SQL & Base de datos administrada\\
Base de datos de documentos & Firestore & NoSQL administrada\\
Ver registros y métricas & Logging y Monitoring & Operaciones\\
\bottomrule
\end{tabular}
\par\vspace{3mm}
\begin{block}{Idea clave}
\small No hay que saberlo todo al inicio. Con estos siete ya puedes seguir los cursos de la ruta.
\end{block}
\end{frame}

\begin{frame}{Cómputo: ¿VM, Cloud Run o GKE?}
\begin{block}{Una comparación para moverse por la ciudad}
\small Las tres llevan tu aplicación a Internet; cambia \textbf{cuánto manejas tú}.
\end{block}
\begin{columns}[T,onlytextwidth]
\column{.31\textwidth}
\begin{exampleblock}{Compute Engine}
\small \textbf{Tu propio auto.} Eliges el modelo, pones la gasolina y lo llevas al taller. Control total, más trabajo.
\end{exampleblock}
\column{.31\textwidth}
\begin{exampleblock}{Cloud Run}
\small \textbf{Un taxi.} Dices a dónde vas y pagas solo el viaje. Si nadie usa tu app, casi no pagas.
\end{exampleblock}
\column{.31\textwidth}
\begin{exampleblock}{GKE}
\small \textbf{Una flota de buses.} Muchos vehículos coordinados. Potente, pero hay que aprender a administrarla.
\end{exampleblock}
\end{columns}
\par\vspace{2mm}
\begin{alertblock}{Pregunta típica del examen}
\small ``Tienes un contenedor, el tráfico es irregular y quieres pagar poco cuando no hay visitas''. Respuesta: \textbf{Cloud Run}.
\end{alertblock}
\end{frame}

\begin{frame}[fragile]{Almacenamiento: Cloud Storage}
\begin{block}{Qué es}
\small Un lugar para guardar \textbf{archivos} (imágenes, PDF, copias de seguridad) dentro de \textbf{buckets}: carpetas con un nombre único en todo el mundo.
\end{block}
\begin{columns}[T,onlytextwidth]
\column{.48\textwidth}
\begin{exampleblock}{Clases de almacenamiento}
\small
\textbf{Standard:} lo que se usa todo el tiempo.\par\smallskip
\textbf{Nearline} y \textbf{Coldline:} lo que se usa poco.\par\smallskip
\textbf{Archive:} lo que casi nunca se abre. Guardar es muy barato; leer cuesta más.
\end{exampleblock}
\column{.48\textwidth}
\begin{block}{Con \texttt{gcloud}}
\begin{lstlisting}[basicstyle=\ttfamily\footnotesize]
gcloud storage buckets create \
  gs://NOMBRE-UNICO \
  --location=us-central1

gcloud storage cp foto.png \
  gs://NOMBRE-UNICO/
\end{lstlisting}
\end{block}
\end{columns}
\par\vspace{1mm}
\begin{alertblock}{Relación con la Sesión 14}
\small Un bucket también puede servir una página estática, como GitHub Pages, pero dentro de Google Cloud.
\end{alertblock}
\end{frame}

\begin{frame}{Costos: el tema que nadie quiere olvidar}
\begin{block}{Idea clave}
\small En la nube pagas por lo que \textbf{usas}, y lo que olvidas encendido \textbf{también se usa}. El Cloud Engineer vigila el gasto.
\end{block}
\begin{columns}[T,onlytextwidth]
\column{.48\textwidth}
\begin{exampleblock}{Para aprender sin gastar}
\small
\textbf{Prueba gratuita:} crédito para cuentas nuevas (hoy, 300 dólares por 90 días).\par\smallskip
\textbf{Nivel gratuito:} algunos servicios tienen un uso mensual gratis, como una VM pequeña \texttt{e2-micro}.\par\smallskip
\textbf{Labs de Google Skills:} usan un proyecto temporal; no se cobra a tu tarjeta.
\end{exampleblock}
\column{.48\textwidth}
\begin{block}{Presupuestos y alertas}
\small En \textbf{Facturación} $\rightarrow$ \textbf{Presupuestos y alertas} defines un monto (por ejemplo, 10 dólares) y Google te avisa por correo al 50\,\%, 90\,\% y 100\,\%.
\end{block}
\begin{alertblock}{Cuidado}
\small El presupuesto \textbf{avisa}, pero \textbf{no apaga} nada solo. Al terminar, detén o borra tus recursos.
\end{alertblock}
\end{columns}
\end{frame}

\begin{frame}{Un plan de estudio para principiantes}
\renewcommand{\arraystretch}{1.35}
\centering\small
\begin{tabular}{@{}p{2.2cm}p{5.4cm}p{5.2cm}@{}}
\toprule
\textbf{Etapa} & \textbf{Qué haces} & \textbf{Dónde}\\
\midrule
Mes 1 & Conceptos: proyectos, regiones, IAM, servicios principales & Curso \emph{Google Cloud Fundamentals: Core Infrastructure}\\
Mes 2 & Práctica guiada: VMs, redes, Storage, Cloud Run & Laboratorios de la ruta Cloud Engineer\\
Mes 3 & Retos sin guía & Insignias de habilidad (\emph{skill badges})\\
Al final & Medir si estás listo & Preguntas de ejemplo oficiales y la guía del examen\\
\bottomrule
\end{tabular}
\par\vspace{4mm}
\begin{columns}[T,onlytextwidth]
\column{.48\textwidth}
\begin{exampleblock}{Hábito que funciona}
\small Un laboratorio corto al día vale más que una maratón el fin de semana.
\end{exampleblock}
\column{.48\textwidth}
\begin{alertblock}{No te apures}
\small Primero entiende; el examen puede esperar. Sin práctica real, la certificación se olvida.
\end{alertblock}
\end{columns}
\end{frame}

% ------------------------------------------------------ Google Skills
\begin{frame}[plain]
\vfill
\begin{center}
{\fontsize{29}{34}\selectfont\bfseries Google Skills: primeros pasos}\par
\vspace{2mm}
{\small Dos laboratorios cortos de la ruta}\par
\vspace{4mm}
{\color{utpRojo}\rule{8.5cm}{1.2pt}}
\end{center}
\vfill
\end{frame}

\begin{frame}[fragile]{Laboratorios de hoy}
\begin{columns}[T,onlytextwidth]
\column{.48\textwidth}
\begin{exampleblock}{Lab 1: Cloud Shell y gcloud}
\small Busca en \texttt{skills.google}: \emph{Getting Started with Cloud Shell and gcloud} (GSP002).\par\smallskip
Configuras región y zona, creas una VM con \texttt{gcloud} y te conectas por SSH.
\end{exampleblock}
\begin{block}{Bloques del examen}
\small 1 (preparar el entorno) y 2 (implementar).
\end{block}
\column{.48\textwidth}
\begin{exampleblock}{Lab 2: Cloud Storage}
\small Busca: \emph{Cloud Storage: Qwik Start - Cloud Console} (GSP073).\par\smallskip
Creas un bucket, subes una imagen y la haces pública.
\end{exampleblock}
\begin{block}{Bloques del examen}
\small 2 (implementar) y 4 (acceso: quién puede ver el archivo).
\end{block}
\end{columns}
\par\vspace{1mm}
\begin{alertblock}{Igual que en la Sesión 14}
\small Abre la consola en \textbf{incógnito} con el usuario temporal y presiona \textbf{Check my progress} en cada tarea.
\end{alertblock}
\end{frame}

\begin{frame}{Práctica: tu ruta Cloud Engineer \textnormal{\small(30 min)}}
\begin{block}{Pasos}
\small
1. Haz el \textbf{Lab 1} (Cloud Shell y gcloud). Anota la región y la zona que usaste.\par\smallskip
2. Haz el \textbf{Lab 2} (Cloud Storage). Abre la imagen pública desde otra pestaña.\par\smallskip
3. Abre \texttt{Sesion15/apps/ruta/index.html}: responde el repaso y marca lo que ya sabes.\par\smallskip
4. Busca en Google Skills la \textbf{ruta Cloud Engineer} y anota los tres primeros cursos o insignias.\par\smallskip
5. Escribe a qué \textbf{bloque del examen} (1 al 4) corresponde cada lab.
\end{block}
\begin{columns}[T,onlytextwidth]
\column{.48\textwidth}
\begin{exampleblock}{Si el bucket no se crea}
\small El nombre ya existe en el mundo: agrega tu código de alumno o números.
\end{exampleblock}
\column{.48\textwidth}
\begin{alertblock}{Entregable}
\small Captura de los dos labs completados y tu lista de los tres primeros pasos de la ruta.
\end{alertblock}
\end{columns}
\end{frame}

\begin{frame}{Cierre: preguntas para comprobar}
\begin{block}{Responde con tus palabras}
\small
1. ¿Qué diferencia hay entre un curso de Google Skills y la certificación Associate Cloud Engineer?\par\medskip
2. Nombra los 4 bloques del examen y da un ejemplo del curso para cada uno.\par\medskip
3. ¿Por qué el \textbf{proyecto} es tan importante en Google Cloud?\par\medskip
4. ¿Qué diferencia hay entre una región y una zona? ¿Una VM vive en cuál?\par\medskip
5. Tienes un contenedor con pocas visitas y quieres pagar poco. ¿Compute Engine, Cloud Run o GKE?\par\medskip
6. Pusiste un presupuesto de 10 dólares. ¿Google apaga tus VMs al llegar a ese monto?
\end{block}
\begin{alertblock}{Lo que viene}
\small Seguimos avanzando en la ruta: redes, contenedores y despliegues en Google Cloud.
\end{alertblock}
\end{frame}

\begin{frame}[plain]
\begin{tikzpicture}[remember picture,overlay]
  \node[anchor=center,inner sep=0pt] at (current page.center)
    {\includegraphics[width=\paperwidth,height=\paperheight]{gracias-final.pdf}};
\end{tikzpicture}
\end{frame}

\begin{frame}{Referencias}
\scriptsize
\begin{itemize}\setlength\itemsep{3pt}
  \item Google Cloud: \emph{Associate Cloud Engineer} (página de la certificación y guía del examen), \texttt{cloud.google.com/learn/certification/cloud-engineer}.
  \item Google Cloud: \emph{Google Cloud certifications}, \texttt{cloud.google.com/learn/certification}.
  \item Google Skills: ruta de aprendizaje \emph{Cloud Engineer} y catálogo de laboratorios, \texttt{skills.google}.
  \item Google Cloud: \emph{Resource hierarchy}, \texttt{cloud.google.com/resource-manager/docs}.
  \item Google Cloud: \emph{Geography and regions} y \emph{Regions and zones}, \texttt{cloud.google.com/compute/docs/regions-zones}.
  \item Google Cloud: \emph{gcloud CLI overview} y \emph{Cloud Shell}, \texttt{cloud.google.com/sdk/docs} y \texttt{cloud.google.com/shell/docs}.
  \item Google Cloud: \emph{Storage classes}, \texttt{cloud.google.com/storage/docs}.
  \item Google Cloud: \emph{Free Program} y \emph{Create, edit, or delete budgets and budget alerts}, \texttt{cloud.google.com/free} y \texttt{cloud.google.com/billing/docs}.
  \item Material de repaso: \texttt{Sesion14/slides/slides.tex} y \texttt{Sesion14/respuestas.md}.
  \item Las respuestas de todas las preguntas de esta sesión están en \texttt{respuestas.md}.
\end{itemize}
\end{frame}

\end{document}
